% !TEX root = ../Main.tex
\chapter{分式不等式}

分式不等式的未知数在分母上。\textbf{不能直接两边乘分母}——分母的正负不知道，乘上去不等号方向就定不下来。

\section{利用完全平方数的非负性}

一个干净的办法：两边同乘分母的\textbf{平方} $g^2$。因为 $g^2\ge 0$（且只要 $g\ne 0$ 就 $g^2>0$），乘上去不会翻转不等号，分式就化成了多项式。

\ex{解 $\dfrac{x-1}{x+1}<0$。}

\sol{
    两边同乘 $(x+1)^2$（严格不等式本身已要求 $x+1\ne 0$）：
    \[
    \frac{x-1}{x+1}<0
    \ \Longleftrightarrow\
    \frac{x-1}{x+1}\,(x+1)^2<0
    \ \Longleftrightarrow\
    (x-1)(x+1)<0.
    \]
    两根 $-1,1$，开口向上，"小于取中间"：
    \[
    \text{解集}=(-1,1).
    \]
}

\ex{解 $\dfrac{x+3}{x-2}\le 0$。}

\sol{
    同乘 $(x-2)^2$，但分母不能为零，要单独排除 $x=2$：
    \[
    \frac{x+3}{x-2}\le 0
    \ \Longleftrightarrow\
    (x+3)(x-2)\le 0\ \text{且}\ x\ne 2.
    \]
    $(x+3)(x-2)\le 0$ 取中间 $-3\le x\le 2$，再去掉 $x=2$：
    \[
    \text{解集}=[-3,2).
    \]
}

\section{右边不是 0：先通分}

如果不等式右边不是 $0$，先\textbf{移项通分}，把一侧化成 $0$、另一侧化成一个分式，再按上面处理。

\ex{解 $\dfrac{x-5}{x+4}\ge 1$。}

\sol{
    移项通分，使右边为 $0$：
    \[
    \frac{x-5}{x+4}-1\ge 0
    \ \Longleftrightarrow\
    \frac{(x-5)-(x+4)}{x+4}\ge 0
    \ \Longleftrightarrow\
    \frac{-9}{x+4}\ge 0.
    \]
    分子 $-9<0$ 是常数，要让分式 $\ge 0$，只能让分母 $x+4<0$（负除以负为正；$x+4$ 不能为零）：
    \[
    \text{解集}=(-\infty,-4).
    \]
}

\rmkb{
    \textbf{方法总结}
    \begin{enumerate}
        \item 先通分，令一侧为 $0$；
        \item 再消去分母（同乘分母的平方），并\textbf{注意分母本身不能为 $0$}。
    \end{enumerate}
}
